% This is the ADASS_template.tex LaTeX file, 24 Oct 2025.
% It is based on the ASP general author template file, but modified to reflect the specific
% requirements of the ADASS proceedings.

% This is the aspauthor.tex LaTeX file
% Copyright 2021, Astronomical Society of the Pacific Conference Series
% Revision:  29 November 2021

% To compile, at the command line positioned at this folder, type:
% pdflatex ADASS_template
% pdflatex ADASS_template
% This will create a file called ADASS_template.pdf

\documentclass[11pt,twoside]{article}

% Do NOT use ANY packages other than asp2023. 
\usepackage{asp2023}

\aspSuppressVolSlug
\resetcounters

% References must all use BibTeX entries in a .bibfile.
% References must be cited in the text using \citet{} or \citep{}.
% Do not use \cite{}.
% See ManuscriptInstructions.pdf for more details
\bibliographystyle{asp2021}

% The ``markboth'' line sets up the running heads for the paper.
% 1 author: "Surname"
% 2 authors: "Surname1 and Surname2"
% 3 authors: "Surname1, Surname2, and Surname3"
% >3 authors: "Surname1 et al."
% Replace ``Short Title'' with the actual paper title, shortened if necessary.
% Use mixed case type for the shortened title
% Ensure shortened title does not cause an overfull hbox LaTeX error
% See ASPmanual2010.pdf 2.1.4  and ManuscriptInstructions.pdf for more details
\markboth{Molinaro, Donaldson, and the IVOA TCG}{IVOA TCG status: 2026}

\begin{document}

\title{Status of the IVOA Technical Standards and Processes as of 2026}

% Note the position of the comma between the author name and the 
% affiliation number.
% Authors surnames should come after first names or initials, eg John Smith, or J. Smith.
% Author names should be separated by commas.
% The final author should be preceded by "and".
% Affiliations should not be repeated across multiple \affil commands. If several
% authors share an affiliation this should be in a single \affil which can then
% be referenced for several author names. If only one affiliation, no footnotes are needed.
% See ManuscriptInstructions.pdf and ASP's manual2010.pdf 3.1.4 for more details
\author{Marco~Molinaro,$^1$ Tom~Donaldson,$^2$ and the IVOA TCG}
\affil{$^1$Institution Name, Institution City, State/Province, Country; \email{AuthorEmail@email.edu}}
\affil{$^2$Institution Name, Institution City, State/Province, Country}

% This section is for ADS Processing.  There must be one line per author. paperauthor has 9 arguments.
\paperauthor{Sample~Author1}{Author1Email@email.edu}{ORCID_Or_Blank}{Author1 Institution}{Author1 Department}{City}{State/Province}{Postal Code}{Country}
\paperauthor{Sample~Author2}{Author2Email@email.edu}{ORCID_Or_Blank}{Author2 Institution}{Author2 Department}{City}{State/Province}{Postal Code}{Country}
\paperauthor{Sample~Author3}{Author3Email@email.edu}{ORCID_Or_Blank}{Author3 Institution}{Author3 Department}{City}{State/Province}{Postal Code}{Country}

% There should be one \aindex line (commented out) for each author. These are used to
% build up the author index for the Proceedings. The surname must come first, followed by
% initials. Note the use of ~ before each initial to control spacing.
% The \author entries (see above) have surname last. These \aindex entries have
% surname first.
% The Aindex.py command willl create them for you after you have constructed the \author
% The first entry should be the first author, for bold-facing the author index page-reference

%\aindex{Molinaro,~M.}
%\aindex{Donaldson,~T.}
%\aindex{IVOA TCG}


\begin{abstract}
In recent years the Technical Coordination Group (TCG) of the International Virtual Observatory Alliance (IVOA) has
continued the activities of developing community open standards targeted at establishing an interoperable scenario in the
general astrophysical research domain.
Alongside the continuous work of revising or defining standards, activity has focused on requirements from research sub-
domains that are re-emerging or have previously received less attention (such as time-domain astronomy or high-energy
astrophysics), on technical challenges arising from new information system technologies (like OpenAPI), and on new data
formats and distribution paradigms (including increased use of tessellation and distributed solutions). The TCG has also
discussed organisational aspects, possible improvements, and challenges of the IVOA standardisation process itself.
This contribution provides a brief overview of all of the above and tries to identify the major priorities that the TCG might
have to face in the oncoming years, balancing the available efforts volunteered by member projects and participating
projects against the requirements for new or renewed standards and their uptake.
\end{abstract}

% These lines show examples of subject index entries. At this stage these have to commented
% out, and need to be on separate lines. Eventually, they will be automatically uncommented
% and used to generate entries in the Subject Index at the end of the Proceedings volume.
% Don't leave these in! - replace them with ones relevant to your paper.
%\ssindex{FOOBAR!conference!ADASS 2019}
%\ssindex{FOOBAR!organisations!ASP}

% These lines show examples of ASCL index entries. At this stage these have to commented
% out, and need to be on separate lines. Eventually, they will be automatically uncommented
% and used to generate entries in the ASCL Index at the end of the Proceedings volume.
% The ascl.py command will scan your paper on possible code names.
% Don't leave these in! - replace them with ones relevant to your paper.
%\ooindex{FOOBAR, ascl:1101.010}


\section{The IVOA Technical Coordination Group}

TCG is one part of the IVOA organisation. As sketched in Fig.~\ref{fig:orgchart}, it has connections to the Executive, the
Standards \& Processes, and the Science Priorities committees.

Recent organisational discussions focus on:
\begin{itemize}
  \item updating the processes to establish standards;
  \item improving transitioning to major versions;
  \item adapting newer technologies in producing the standards (like OpenAPI for standard interfaces);
  \item clarifying the role of Interest Groups with respect to Science Priorities.
\end{itemize}

A dedicated tiger team was dedicated to transitioning solutions to better interface specifications, the Standard
\& Processes committee is working on updating the processes.

\articlefigure[width=0.65\textwidth]{ivoa-orgchart.eps}{fig:orgchart}{IVOA organisational structure overview. The TCG is expanded in its Working and Interest Groups components.}

\section{Recent TCG activities}

The most relevant goal of the TCG is to coordinate and produce standards (IVOA Recommendations\footnote{See them listed in the current IVOA Documents' repository: \url{https://www.ivoa.net/documents}}). In the last few years this activity has produced new or updated standards (see Fig.~\ref{fig:newstds} for a quick overview of the standards that have been recently touched), or connected activities, like:
\begin{itemize}
  \item semantic updates;
  \item vocabularies and their maintenance;
  \item connection to external technical standards;
  \item OpenAPI considerations;
  \item .parquet file integration with VO metadata;
  \item enrichment of VOTable responses with provenance information;
  \item solutions to integrate data collections into bibliographic repositories (ADS);
  \item extension of the Observational Core model to different wavelength/energy regimes;
  \item inclusion of heliospheric physics and planetary sciences requirements;
  \item reviving effort in time-domain astronomy;
  \item look-out to external community efforts (ASDF, ADPFG).
\end{itemize}

\articlefigure[width=0.65\textwidth]{recent-stds.eps}{fig:newstds}{IVOA Recommendations updates in the last three years.}

\section{Future activities}

Still, a lot needs to be pursued, also considering the volunteer-based effort at IVOA:
\begin{itemize}
  \item tessellation and new data formats;
  \item pursue OpenAPI as a way to make more robust protocol technical description;
  \item discuss and define the scope of Recommendations and Endorsed Notes;
  \item connect to other communities and improve uptake;
  \item standards updates: ConeSearch, SLAP, VOTable, CAOM, ObsTAP, \ldots;
  \item revive simulated data standardisation efforts;
  \item AI connection to VO (follow ai-interop@ivoa.net).
\end{itemize}

\acknowledgements MM and TD want to acknowledge all the TCG members for the work done. For an up to date view of the active roles in coordination, check the IVOA TCG\footnote{\url{https://wiki.ivoa.net/twiki/bin/view/IVOA/IvoaTCG}} dedicated page.

\bibliography{P602}  % For BibTex

% if we have space left, we might add a conference photograph here. Leave commented for now.
% \bookpartphoto[width=1.0\textwidth]{foobar.eps}{FooBar Photo (Photo: Any Photographer)}

\end{document}
